\SourceLeaf{642}

II. Der Flächenraum $s$ sei von der Linie $u$ begrenzt, $\mathrm du$ ein. Element von $u$; $R,r$ seine Entfernung von $P$ und $p$; $w$ der Winkel zwischen $R$ und $r$, und $v$ der Winkel zwischen $u$ und dem Planum durch $R,r$.

Dann ist der Vibrationszustand in $p$
\[
=\int\frac{\mathrm du.\sin v}{Rr.\sin w}e^{\alpha(R+r)}
\]
Kürzer so:

es seien $x,y,z$ Coordinaten jedes Punktes im Raume und für
\[
\begin{array}{lccc}
P\, .\, .\, .&0,&0,&0\\
p\, .\, .\, .&0,&0,&h
\end{array}
\]
dann ist der Vibrationszustand in $p$
\[
=\int\frac{x\mathrm dy-y\mathrm dx}{xx+yy}.e^{\alpha(R+r)}:h
\]
durch die Randlinie ausgedehnt: oder kürzer, wenn $\dfrac yx=\operatorname{tg}\theta$,
\[
=\int e^{\alpha(R+r)}\mathrm d\theta:h
\]

\SourceLeaf{643}
\section*{BEMERKUNGEN.}
\addcontentsline{toc}{section}{Bemerkungen}

Die hier unter 21 Nummern zusammengestellten bruchstücksweise aufgezeichneten Untersuchungen gehören ziemlich weit auseinanderliegenden Zeiten an. Nr. 1 und 2. befinden sich in einem Tagebuche zwischen den Protocollen von Beobachtungen, die im März, Juni und Juli 1833 über die durch Magnete inducirten Galvanischen Ströme angestellt sind. Nr. 6. 12. 18. 21. stehen auf besondern Blättern und lassen ausser 18., welches ein Datum trägt, keine besondere Zeitbestimmung zu. Die übrigen Nummern mit Ausschluss des Briefes [20.] sind hier in gleicher Reihenfolge wieder gegeben, wie sie sich in einem Handbuche befinden, wo sie aber zahlreiche ganz heterogene Entwickelungen zwischen sich enthalten. Die letzte jener Nummern mit dem Beweise von \textsc{Ampères} Fundamentalsatz ist erst nach 1843 eingetragen, die andern scheinen der Zeit von 1833 bis 1836 anzugehören.

Die verschiedenen Formen, welche hier für das Gesetz der Wechselwirkungen zwischen Galvanischen Stromelementen angenommen werden, ergeben sich alle aus dem besonders in Nr. 20. hervorgehobenen Princip der Umtauschbarkeit des Magnetismus mit galvanischen Strömen. Die in diesem Briefe angedeuteten Untersuchungen von \textsc{Wilhelm Weber} bilden die Vorarbeiten zu der (im Jahre 1846, in der ersten Abhandlung über Electrodynamische Maassbestimmungen, vollendeten) Aufstellung einer Theorie, nach welcher die ganze Wechselwirkung zwischen zwei (mit dem entsprechenden Vorzeichen versehenen $e$, $e'$) electrischen Theilchen, in der gegenseitigen Entfernung $r$, durch
\[
ee'\left(\frac1{rr}-\frac1{cc}\frac1{rr}\frac{\mathrm dr^2}{\mathrm dt^2}+\frac2{cc}\frac1r\frac{\mathrm ddr}{\mathrm dt^2}\right)
\]
gemessen wird und ein positiver Werth dieser Grösse eine Abstossung, ein negativer eine Anziehung bedeutet. In dem Ausdrucke bezeichnet $t$ die Zeit, $c$ eine Geschwindigkeit, welche \textsc{Kohlrausch} und \textsc{Weber} durch Untersuchungen (1855) zur Zurückführung der Stromintensitäts-Messungen auf mechanisches Maass gleich $439450.10^6\,\dfrac{\text{Millimeter}}{\text{Secunde}}$ gefunden haben.

\SourceLeaf{644}

Dem Lehrsatze in Nr. 4 ist eine rein geometrische Einkleidung gegeben; wegen seiner Wichtigkeit für die Theorie der galvanischen Ströme glaubte ich ihm diese Stelle zuweisen zu müssen. Das Integral, durch welches die Anzahl der Umschlingungen der geschlossenen Curve $s$ mit dem System geschlossener Curven $s'$ bestimmt wird, gibt nemlich, wenn statt $\mathrm ds'$ die Elemente aller im Raume vorhandenen geschlossenen galvanischen Ströme gesetzt werden, die algebraische Summe der Intensitäten derjenigen unter diesen Strömen, welche eine von $s$ begrenzte aber im übrigen beliebig bestimmt angenommene Fläche durchdringen. Das Integral selbst ist aber nach dieser Deutung der $\mathrm ds'$ gleich $\dfrac1{4\pi}\int\dfrac{\mathrm dV}{\mathrm ds}\mathrm ds$, wenn $V$ die Potentialfunction für die magnetischen Wirkungen der Ströme $s'$ wie in Nr. 9 bezeichnet. Der Satz bildet also das Analogon zu dem von \textsc{Gauss} in der Abhandlung über die Attraction der Ellipsoide aufgestellten, welcher die innerhalb einer geschlossenen Fläche $(\omega)$ befindliche Masse aus der zur Fläche nach innen gerichteten Normalkräften ihrer Attraction $\left(\dfrac{\mathrm dV}{\mathrm dN}\right)$ durch $\dfrac1{4\pi}\int\dfrac{\mathrm dV}{\mathrm dN}\mathrm d\omega$ bestimmt.

Die Ermittelung des angedeuteten Werthes des obigen Integrals $\dfrac1{4\pi}\allowbreak\int\dfrac{\mathrm dV}{\mathrm ds}\mathrm ds$ ergibt sich z. B. wenn man die Integral-Ausdrucke für die Derivirten von $V$ nach den Coordinaten verwandelt in Integrale, welche sich über irgend beliebig bestimmt angenommene von den einzelnen Stromleitern $s'$ begrenzte Flächen $\omega$ erstrecken. Die dadurch erhaltene Form für die Derivirte von $V$ nach $s$ lässt nach einem im Art. 38 der allgemeinen Theorie des Erdmagnetismus angedeuteten Satze, welcher den Unterschied der Werthe der Potentialfunction für eine auf beiden Seiten mit entgegengesetztem magnetischen Fluidum in geeigneter Weise belegte Fläche und zwar der Werthe an entsprechenden Stellen der beiden Seiten der Fläche angibt, unmittelbar erkennen, dass das gesuchte Integral gleich ist der algebraischen Summe der Intensitäten der Ströme, welche in den Begrenzungslinien der von der Curve $s$ durchsetzten Flächen $\omega'$ sich bewegen.

Die Verwandlung der über eine geschlossene Curve $s$ ausgedehnten Integrale in solche, die sich auf eine von $s$ begrenzte Fläche $\omega$ beziehen, kann mit Hülfe des Satzes ausgeführt werden, dass für irgend welche rechtwinklige gerad- oder krummlinige Coordinaten $\xi,\eta,\zeta$, die also das Quadrat des Längenelements allgemein durch einen Ausdruck von der Form
\[
\xi'\xi'\mathrm d\xi^2+\eta'\eta'\mathrm d\eta^2+\zeta'\zeta'\mathrm d\zeta^2
\]
darstellen, und für beliebige mit ihren ersten Derivirten in den Punkten der Fläche $\omega$ stetig veränderliche Functionen $\lambda,\mu,\nu$ der Coordinaten $\xi,\eta,\zeta$ immer
\[
\begin{aligned}
&\int\left(\lambda\frac{\mathrm d\xi}{\mathrm ds}+\mu\frac{\mathrm d\eta}{\mathrm ds}+\nu\frac{\mathrm d\zeta}{\mathrm ds}\right)\mathrm ds\\
&=\int\left\{\left(\frac{\mathrm d\mu}{\mathrm d\zeta}-\frac{\mathrm d\nu}{\mathrm d\eta}\right)\frac{\xi'}{\eta'\zeta'}\frac{\mathrm d\xi}{\mathrm dn}+\left(\frac{\mathrm d\nu}{\mathrm d\xi}-\frac{\mathrm d\lambda}{\mathrm d\zeta}\right)\frac{\eta'}{\zeta'\xi'}\frac{\mathrm d\eta}{\mathrm dn}\right.\\
&\qquad\left.+\left(\frac{\mathrm d\lambda}{\mathrm d\eta}-\frac{\mathrm d\mu}{\mathrm d\xi}\right)\frac{\zeta'}{\xi'\eta'}\frac{\mathrm d\zeta}{\mathrm dn}\right\}\mathrm d\omega
\end{aligned}
\]
ist, wenn $\xi,\eta,\zeta$ im ersten Integral die Coordinaten eines Punktes des Längenelements $\mathrm ds$, im zweiten $\xi,\eta,\zeta$ die Coordinaten eines Punktes des Flächenelements $\mathrm d\omega$ und $n$ die Normale zu diesem Flächenelement bedeuten. Die positive Richtung der Normale ist so zu wählen, dass, wenn $\mathrm dt$ das erste Element einer von einem Punkte des $\mathrm ds$ zu diesem selbst normal aber in der Fläche $\omega$ liegenden Curve bezeichnet, die positiven Richtungen der $\xi'\mathrm d\xi$, $\eta'\mathrm d\eta$, $\zeta'\mathrm d\zeta$ durch stetige Verschiebung der Lage des Coordinatensystems im Raume der Reihe nach mit den positiven Richtungen der $\mathrm dn$, $\mathrm ds$, $\mathrm dt$ zur Deckung gebracht werden können.

\SourceLeaf{645}

Dieser Satz gibt durch wiederholte Anwendung auch den Beweis von \textsc{Ampères} Fundamentalsatz in der allgemeinen Form, dass unmittelbar die Potentialfunction für die Wechselwirkung zwischen den auf bestimmte Weise mit magnetischem Fluidum belegten Flächen zurückgeführt wird auf die Potentialfunction für die Wechselwirkung zwischen galvanischen Strömen, die nach Lage und Intensität durch jene Flächen und die Magnetisirung bestimmt sind.

Dem in Nr. 9. aufgestellten Beweise für die Gleichheit der Werthe der verschiedenen Ausdrücke für die Potentialfunction $V$ kann man eine symmetrische Form geben, wenn man die Function $R$
\[
\begin{aligned}
={}&xx\operatorname{arctang}\frac{yz}{xr}+yy\operatorname{arctang}\frac{zx}{yr}+zz\operatorname{arctang}\frac{xy}{zx}\\
&+2yzi\operatorname{arctang}\frac{xi}r+2zxi\operatorname{arctang}\frac{yi}r+2xyi\operatorname{arctang}\frac{zi}r
\end{aligned}
\]
worin $i$ statt $\BareRoot{-1}$ gesetzt ist, einführt und berücksichtigt, dass die Gleichungen
\[
\begin{array}{ll}
\dfrac{\mathrm ddR}{\mathrm dx^2}=2\operatorname{arctang}\dfrac{yz}{xr},&\dfrac{\mathrm ddR}{\mathrm dy\mathrm dz}=2i\operatorname{arctang}\dfrac{xi}r\\[5pt]
\dfrac{\mathrm ddR}{\mathrm dy^2}=2\operatorname{arctang}\dfrac{zx}{yr},&\dfrac{\mathrm ddR}{\mathrm dz\mathrm dx}=2i\operatorname{arctang}\dfrac{yi}r\\[5pt]
\dfrac{\mathrm ddR}{\mathrm dz^2}=2\operatorname{arctang}\dfrac{xy}{zr}.&\dfrac{\mathrm ddR}{\mathrm dx\mathrm dy}=2i\operatorname{arctang}\dfrac{zi}r
\end{array}
\]
\[
\frac{\mathrm ddR}{\mathrm dx^2}+\frac{\mathrm ddR}{\mathrm dy^2}+\frac{\mathrm ddR}{\mathrm dz^2}=(2m+1)\pi,\qquad\frac{\mathrm d^3R}{\mathrm dx\mathrm dy\mathrm dz}=-2\frac1r
\]
Statt haben. Durch die Derivirten der Function $R(x,y,z)$ können in endlicher Form auch die Potentialfunctionen für die magnetische Wirkung solcher galvanischer Ströme dargestellt werden, deren Leiter aus geradlinigen den Axen eines rechtwinkligen Coordinatensystems parallelen linearen Theilen bestehen.

Die von \textsc{Gauss} bei der Bestimmung einer Abtheilung eines getheilten Maassstabes in Theilen des ganzen Maassstabes angewandte Anordnung der Längen-Comparirungen danken wir der Aufzeichnung, die sich der Herr Geh. Hofrath \textsc{Weber} im Jahre 1839 oder 1840 gemacht hat.

Die über die Beugungserscheinungen angestellten theoretischen Untersuchungen sind wahrscheinlich durch das von \textsc{F. M. Schwerd} im Jahre 1835 herausgegebene diesen Gegenstand betreffende Werk veranlasst. Die beiden für die Wirkung eines leuchtenden Punkts $P$ auf einen Punkt $p$ aufgestellten allgemeinen Formeln sind nicht identisch; die allgemeine Verwandlung solcher Flächen-Integrale, deren Elemente von der Lage der durch einen Punkt des zugehörigen Flächentheils $\mathrm ds$ und durch die Punkte $P$ und $p$ gehenden Ebene nicht abhangen, in solche Curven-Integrale, deren Elemente ebenfalls von der Lage der durch einen Punkt des zugehörigen Theilchens der Begrenzungslinie $u$ und durch die Punkte $P$ und $p$ gehenden Ebene nicht abhangen, deren Differentiale aber eine Änderung allein des Winkels $\theta$ bedeuten, welchen jene Ebene mit einer durch $P$ und $p$ gelegten festen Ebene einschliesst, ergibt sich aus der Gleichung
\SourceLeaf{646}
\[
\begin{aligned}
\int Q\mathrm d\theta=\int\frac{h.Q.\sin v.\mathrm du}{rR\sin w}
={}&\int\frac h{rR}.\frac1{2\sin\frac12w^2}.\frac{\mathrm dQ}{\mathrm d(R-r)}\frac{\mathrm d(R+r)}{\mathrm d\rho}.\mathrm ds\\
&-\int\frac h{rR}.\frac1{2\cos\frac12w^2}.\frac{\mathrm dQ}{\mathrm d(R+r)}.\frac{\mathrm d(R-r)}{\mathrm d\rho}\,\mathrm ds
\end{aligned}
\]
die einen speciellen Fall des in der vorhergehenden Bemerkung erwähnten Satzes bildet, wenn nemlich $Q$ eine von $R$ und $r$ allein abhängige mit ihren nach $R+r$ und $R-r$ genommenen partiellen Derivirten für die Punkte der Fläche $s$ stetig veränderliche Grösse bedeutet, ferner $w$ genauer als im Text dahin bestimmt ist, dass es den Winkel zwischen $R$ und $r$ bezeichnet, den die mit den Richtungen des fortschreitenden Lichtstrahls übereinstimmend angenommenen positiven Richtungen jener Linien einschliessen.

\begin{flushright}\textsc{Schering.}\end{flushright}

\clearpage
\SourceLeaf{647}
\section*{INHALT.}
\addcontentsline{toc}{section}{Inhalt}
\begin{center}\bfseries GAUSS WERKE BAND V.\par MATHEMATISCHE PHYSIK.\end{center}
\begingroup\fontsize{9}{11}\selectfont\raggedright
\par\Needspace{5\baselineskip}\medskip\noindent\textit{Abhandlungen}\par\smallskip
\noindent\hangindent=1em Theoria attractionis corporum sphaeroidicorum ellipticorum homogeneorum \unskip\hspace{1em plus 1fill}\mbox{1813 März . .\quad Seite 1}\par\smallskip
\noindent\hangindent=1em Uber ein neues allgemeines Grundgesetz der Mechanik \unskip\hspace{1em plus 1fill}\mbox{1829\quad — 23}\par\smallskip
\noindent\hangindent=1em Principia generalia theoriae figurae fluidorum in statu aequilibrii . . .\unskip\hspace{1em plus 1fill}\mbox{1829 Sept. . .\quad — 29}\par\smallskip
\noindent\hangindent=1em Intensitas vis magneticae terrestris ad mensuram absolutam revocata . . .\unskip\hspace{1em plus 1fill}\mbox{1832 Dec. . .\quad — 79}\par\smallskip
\noindent\hangindent=1em Allgemeine Theorie des Erdmagnetismus . . . .\unskip\hspace{1em plus 1fill}\mbox{1838\quad 119}\par\smallskip
\noindent\hangindent=1em Allgemeine Lehrsätze in Beziehung auf die im verkehrten Verhältnisse des Quadrats der Entfernung wirkenden Anziehungs- und Abstossungs-Kräfte \unskip\hspace{1em plus 1fill}\mbox{1839\quad — 195}\par\smallskip
\noindent\hangindent=1em Dioptrische Untersuchungen . . . . . . . . . . . . . .\unskip\hspace{1em plus 1fill}\mbox{1840 Dec. . .\quad — 243}\par\smallskip
\par\Needspace{5\baselineskip}\medskip\noindent\textit{Anzeigen eigner Abhandlungen.}\par\smallskip
\noindent\hangindent=1em Theoria attractionis corporum sphaeroidicorum ellipticorum homogeneorum .\unskip\hspace{1em plus 1fill}\mbox{1813 April . .\quad — 279}\par\smallskip
\noindent\hangindent=1em Principia generalia theoriae figurae fluidorum in statu aequilibrii . . . .\unskip\hspace{1em plus 1fill}\mbox{1829 Oct. . .\quad — 287}\par\smallskip
\noindent\hangindent=1em Intensitas vis magneticae terrestris ad mensuram absolutam revocata . . .\unskip\hspace{1em plus 1fill}\mbox{1832 Dec. . .\quad — 293}\par\smallskip
\noindent\hangindent=1em Allgemeine Lehrsätze in Beziehung auf die im verkehrten Verhältnisse u. s. w. \unskip\hspace{1em plus 1fill}\mbox{1840 März .\quad — 305}\par\smallskip
\noindent\hangindent=1em Dioptrische Untersuchungen . . . . . . . . . . . . . . .\unskip\hspace{1em plus 1fill}\mbox{1841 Jan. . .\quad — 309}\par\smallskip
\par\Needspace{5\baselineskip}\medskip\noindent\textit{Verschiedene Aufsätze über Magnetismus.}\par\smallskip
\noindent\hangindent=1em Erdmagnetismus und Magnetometer . . . . . . . . . .\unskip\hspace{1em plus 1fill}\mbox{1836\quad — 315}\par\smallskip
\noindent\hangindent=1em Einleitung für die Zeitschrift: Resultate u. s. w. . . . . . . . .\unskip\hspace{1em plus 1fill}\mbox{1836\quad — 345}\par\smallskip
\noindent\hangindent=1em Ein neues Hülfsmittel für die magnetischen Beobachtungen . . . . . .\unskip\hspace{1em plus 1fill}\mbox{1837 Oct\quad — 352}\par\smallskip
\noindent\hangindent=1em Uber ein neues, zunächst zur unmittelbaren Beobachtung der Veränderungen in der Intensität des horizontalen Theils des Erdmagnetismus bestimmten Instruments . . . . . . . . . . . . . . . . .\unskip\hspace{1em plus 1fill}\mbox{1837\quad — 357}\par\smallskip
\noindent\hangindent=1em Anleitung zur Bestimmung der Schwingungsdauer einer Magnetnadel . .\unskip\hspace{1em plus 1fill}\mbox{1837\quad — 374}\par\smallskip
\noindent\hangindent=1em Über ein Mittel die Beobachtung von Ablenkungen zu erleichtern . .\unskip\hspace{1em plus 1fill}\mbox{1839\quad — 395}\par\smallskip
\noindent\hangindent=1em Zur Bestimmung der Constanten des Bifilarmagnetometers . . . . .\unskip\hspace{1em plus 1fill}\mbox{1840\quad — 404}\par\smallskip
\noindent\hangindent=1em Vorschriften zur Bestimmung der magnetischen Wirkung, welche ein Magnetstab in der Ferne ausübt . . . . . . . . . . . . . .\unskip\hspace{1em plus 1fill}\mbox{1840\quad — 427}\par\smallskip
\SourceLeaf{648}
\noindent\hangindent=1em Über die Anwendung des Magnetometers zur Bestimmung der absoluten Declination . . . . . . . . . . . . . . . . . . . .\unskip\hspace{1em plus 1fill}\mbox{1841\quad Seite 436}\par\smallskip
\noindent\hangindent=1em Beobachtungen der magnetischen Inclination in Göttingen . . . . . .\unskip\hspace{1em plus 1fill}\mbox{1841\quad — 444}\par\smallskip
\par\Needspace{5\baselineskip}\medskip\noindent\textit{Aufsätze über verschiedene Gegenstände der mathematischen Physik.}\par\smallskip
\noindent\hangindent=1em Fundamentalgleichungen für die Bewegung schwerer Körper auf der rotirenden Erde . . . . . . . . . . . . . . . . . . . .\unskip\hspace{1em plus 1fill}\mbox{1803\quad — 495}\par\smallskip
\noindent\hangindent=1em Über die achromatischen Doppelobjective besonders in Rücksicht der vollkommenern Aufhebung der Farbenzerstreuung . . . . . . . .\unskip\hspace{1em plus 1fill}\mbox{1817 Dec. . .\quad — 504}\par\smallskip
\noindent\hangindent=1em Brief an \textsc{Brandes} über denselben Gegenstand . . . . . . . . . .\unskip\hspace{1em plus 1fill}\mbox{1831\quad — 509}\par\smallskip
\noindent\hangindent=1em Berichtigung der Stellung der Schneiden einer Waage . . . . . . .\unskip\hspace{1em plus 1fill}\mbox{1837 März . .\quad — 511}\par\smallskip
\par\Needspace{5\baselineskip}\medskip\noindent\textit{Physicalische Beobachtungen.}\par\smallskip
\noindent\hangindent=1em Nordlicht am 7. Januar 1831 . . . . . . . . . . . . . . .\unskip\hspace{1em plus 1fill}\mbox{1831 Febr. . .\quad — 517}\par\smallskip
\noindent\hangindent=1em Magnetisches Observatorium in Göttingen . . . . . . . . . . .\unskip\hspace{1em plus 1fill}\mbox{1834 Aug. . .\quad — 519}\par\smallskip
\noindent\hangindent=1em Beobachtungen der magnetischen Variation in Göttingen und Leipzig . .\unskip\hspace{1em plus 1fill}\mbox{1834 Oct. . .\quad — 523}\par\smallskip
\noindent\hangindent=1em Magnetisches Observatorium in Göttingen . . . . . . . . . . .\unskip\hspace{1em plus 1fill}\mbox{1835 März . .\quad — 528}\par\smallskip
\noindent\hangindent=1em Beobachtungen der magnetischen Variation in Copenhagen und Mailand .\unskip\hspace{1em plus 1fill}\mbox{1835 März . .\quad — 537}\par\smallskip
\noindent\hangindent=1em Magnetisches Observatorium in Göttingen . . . . . . . . . . .\unskip\hspace{1em plus 1fill}\mbox{1836 Juni . .\quad — 540}\par\smallskip
\noindent\hangindent=1em Das in den Beobachtungsterminen anzuwendende Verfahren . . . . .\unskip\hspace{1em plus 1fill}\mbox{1836\quad — 541}\par\smallskip
\noindent\hangindent=1em Auszug aus dreijährigen täglichen Beobachtungen der magnetischen Declination zu Göttingen . . . . . . . . . . . . . . . . .\unskip\hspace{1em plus 1fill}\mbox{1836\quad — 556}\par\smallskip
\noindent\hangindent=1em Erläuterungen zu den Terminszeichnungen und den Beobachtungszahlen .\unskip\hspace{1em plus 1fill}\mbox{1836\quad — 568}\par\smallskip
\noindent\hangindent=1em Erläuterungen zu den Terminszeichnungen und den Beobachtungszahlen .\unskip\hspace{1em plus 1fill}\mbox{1837\quad — 576}\par\smallskip
\noindent\hangindent=1em Der magnetische Südpol der Erde . . . . . . . . . . . . .\unskip\hspace{1em plus 1fill}\mbox{1841\quad — 580}\par\smallskip
\par\Needspace{5\baselineskip}\medskip\noindent\textit{Anzeigen nicht eigner Schriften.}\par\smallskip
\noindent\hangindent=1em \textsc{Benzenberg.} Über die \textsc{Dalton}sche Theorie . . . . . . . . . .\unskip\hspace{1em plus 1fill}\mbox{1830 Dec. . .\quad — 583}\par\smallskip
\noindent\hangindent=1em \textsc{Fischer.} Künstliche Magnete . . . . . . . . . . . . . . .\unskip\hspace{1em plus 1fill}\mbox{1832 Sept. .\quad — 591}\par\smallskip
\noindent\hangindent=1em Resultate aus den Beobachtungen des magnetischen Vereins . . . . .\unskip\hspace{1em plus 1fill}\mbox{1837 Juni . .\quad — 595}\par\smallskip
\noindent\hangindent=1em \textsc{Gerling.} Einrichtung des mathematisch-physicalischen Instituts . . .\unskip\hspace{1em plus 1fill}\mbox{1842 April .\quad — 596}\par\smallskip
\par\Needspace{5\baselineskip}\medskip\noindent\textit{Nachlass.}\par\smallskip
\noindent\hangindent=1em Zur Electrodynamik . . . . . . . . . . . . . . . . . . . . . . . .\unskip\hspace{1em plus 1fill}\mbox{\quad — 601}\par\smallskip
\noindent\hangindent=1em Über Kugelfunctionen . . . . . . . . . . . . . . . . . . . . . . . .\unskip\hspace{1em plus 1fill}\mbox{\quad — 630}\par\smallskip
\noindent\hangindent=1em Zum Gebrauch des Comparators . . . . . . . . . . . . . . . . . . . .\unskip\hspace{1em plus 1fill}\mbox{\quad — 632}\par\smallskip
\noindent\hangindent=1em Allgemeine Formeln für die Wirkung eines leuchtenden Punkts $P$ auf einen Punkt $p$ . .\unskip\hspace{1em plus 1fill}\mbox{\quad — 635}\par\smallskip
\noindent\hangindent=1em Bemerkungen . . . . . . . . . . . . . . . . . . . . . . . . . . . .\unskip\hspace{1em plus 1fill}\mbox{\quad — 637}\par\smallskip
\par\medskip\noindent Steindrucktafel zur Theorie des Erdmagnetismus Seite 176.
\endgroup

\bigskip
\begin{center}
GÖTTINGEN,\par
\small GEDRUCKT IN DER DIETERICHSCHEN UNIVERSITÄTS-DRUCKEREI\par
W. FR. KAESTNER.
\end{center}
